% ECONOMICS_PROBLEM: {"id": "D31-357", "title": "The true global wealth distribution", "jel_code": "D31", "source_id": "OP-0361"}
% !TeX program = xelatex
\documentclass[11pt,a4paper]{article}
\usepackage[margin=23mm]{geometry}
\usepackage{fontspec}
\setmainfont{Latin Modern Roman}
\usepackage{amsmath,amssymb,mathtools}
\usepackage{xcolor,enumitem,booktabs,longtable,array,xurl,hyperref}
\definecolor{muted}{HTML}{53636C}
\hypersetup{colorlinks=true,urlcolor=blue,linkcolor=blue}
\setlength{\parindent}{0pt}
\setlength{\parskip}{5pt}
\newcommand{\R}{\mathbb R}
\newcommand{\E}{\mathbb E}
\newcommand{\N}{\mathbb N}
\newcommand{\ind}{\mathbf 1}
\newcommand{\Var}{\operatorname{Var}}
\newcommand{\Cov}{\operatorname{Cov}}
\newcommand{\supp}{\operatorname{supp}}
\DeclareMathOperator*{\argmin}{arg\,min}
\DeclareMathOperator*{\argmax}{arg\,max}
\newcommand{\entrymeta}[3]{{\sffamily\small\color{muted}\textbf{\texttt{#1}}\quad|\quad #2\par #3\par}}
\newcommand{\Problemref}[1]{\texttt{#1}}
\newcommand{\JELref}[2]{\texttt{#2} (JEL \texttt{#1})}
\begin{document}
\section*{The true global wealth distribution}
\textbf{Problem ID:} \texttt{D31-357}\quad \textbf{JEL:} \texttt{D31}\par
% BEGIN ECONOMICS STATEMENT
\label{problem:OP-0361}
{\sffamily\small\color{muted}Original subject group: Inequality and economic mobility\par}
\label{definitions:problem:30007643}
\textbf{Audit classification:} Published research agenda. Author-hosted published PDF checked. Worldwide hidden ownership motivates the measurement agenda; the quantile-based sharp-bound formulation is editorial.
\subsection*{Definitions and precise research target}
\textbf{Complete editorial problem.} Fix a date and define the global population as all resident adults, assigning each person to one country of residence. For person \(i\), let \(A_i\) be the market value of assets beneficially owned worldwide and \(D_i\) debts; net wealth is \(W_i=A_i-D_i\). Beneficial ownership rules must specify trusts, foundations, private businesses, and jointly owned property. Let \(F_W\) be the resulting world distribution after converting values with a declared exchange-rate or purchasing-power rule. Given household surveys, registries, estate records, international financial positions, and enforcement disclosures, recover the top-share function \(S_p=\int_{1-p}^{1}F_W^{-1}(u)\,du/E[W]\), where the generalized inverse uses the declared population weights, finite absolute mean wealth is assumed, and the denominator is positive. The quantile integral allocates boundary ties fractionally. Construct identified intervals allowing offshore assets, rich-household nonresponse, uncertain private-business values, and missing ownership links. Reconcile domestic and foreign asset aggregates without double-counting custodial holdings or corporations. Determine which additional measurements narrow the intervals most. The target is a date-specific distribution with quantified uncertainty, not a proof that an administrative total uniquely identifies ownership. Earlier global estimates are substantive evidence but do not eliminate the remaining coverage and valuation problem.

\textbf{Definitions and admissible scope.} Let \(\mu\) be the uniform global-adult population measure at a fixed date, or documented sampling weights summing to one. Wealth is beneficially owned asset value minus liabilities, with trusts and joint ownership allocated once. Use \(S_p=\int_{1-p}^1F_W^{-1}(u)\,du/E[W]\), assuming \(E|W|<\infty\) and \(E[W]>0\), instead of including everyone at a tied cutoff. Fractions at quantile ties are allocated by the integral. Negative wealth is allowed and a top share may exceed one. Exchange-rate valuation is the baseline; a purchasing-power conversion produces a different declared concept. An admissible distribution must satisfy ownership-link, valuation, balance-sheet, and missing-asset constraints. Sharpness means each bound is attainable, or approached, within that constraint set. Fix \(p\in(0,1)\). Let \(\mathcal F_W(P_D)\) be the nonempty class of wealth distributions and ownership assignments generating the observed source law under declared selection, linkage, valuation, and missing-asset constraints. The output is the image \(\{S_p(F):F\in\mathcal F_W(P_D)\}\) and its sharp lower/upper endpoints for each requested \(p\). The ownership system treats each liability and each beneficial asset claim exactly once and uses the same date and currency across sources. A measurement-design extension chooses an additional observation law at a stated cost and minimizes a declared criterion such as expected endpoint distance; this is conditional on a prior or minimax class supplied as input.
\subsection*{Variants and assumption changes}
\begin{enumerate}
\item Offshore ownership version (source-motivated): treat missing offshore asset totals as known intervals and their beneficial-owner distributions as restricted families. Bound top shares under those joint constraints and avoid counting custodians as owners.
\item Valuation and measurement version (editorial): allow private-business prices and rich-household response probabilities to vary over stated bounds. Optimize the expected reduction in interval width from a fixed-cost audit sample or ownership registry update.
\end{enumerate}
\subsection*{References and source audit}
\textbf{Support and access.} Author-hosted published PDF checked. Worldwide hidden ownership motivates the measurement agenda; the quantile-based sharp-bound formulation is editorial. \textbf{Locator:} Section 6, Conclusion, p. 134.
\begin{enumerate}\raggedright
\item\hypertarget{definitions:source:30007643:1}{} Gabriel Zucman. 2019. \emph{Global Wealth Inequality}. Section 6, Conclusion, printed p. 134.
\url{https://gabriel-zucman.eu/files/Zucman2019.pdf}.
DOI: \url{https://doi.org/10.1146/annurev-economics-080218-025852}.
\end{enumerate}

\subsection*{Common conventions: Source mathematical conventions}

These conventions apply to each entry unless explicitly replaced.

\textbf{Models and identification.} A primitive model \(m\) specifies all probability laws, feasible choices, information, equilibrium and measurement rules used by its target. The admissible class \(\mathcal M\) is fixed independently of outcomes. If \(P_O(m)\) is its observable probability law and \(\tau(m)\) the target, its identified set at observable law \(P\) is
\[
 \mathcal I_\tau(P)=\{\tau(m):m\in\mathcal M,\ P_O(m)=P\}.
\]
Point identification means that this set is a singleton. An empty set signals incompatibility of data and assumptions. Coordinate bounds need not describe the joint set. Bounds are sharp only if every retained target is attainable or arbitrarily approachable by compatible models. Finite bounds require explicit restrictions on unobserved quantities; unrestricted confounding, hidden wealth, extrapolation or arbitrary equilibrium selection can yield uninformative sets.

\textbf{Causal interventions.} A potential outcome \(Y_i(p)\) is the outcome under a complete policy regime \(p\), including timing, financing, information and market spillovers. The target population and units are fixed before treatment. With interference, use the complete assignment vector or a justified exposure mapping; individual no-interference notation alone cannot identify an economy-wide intervention. A difference of mean potential outcomes does not require the cross-world joint distribution. Conditioning on post-treatment migration, survival, enrollment or employment changes the estimand and needs separate justification.

\textbf{Statistical statements.} An estimator, confidence set, forecast or test specifies a sampling law, model class and loss/coverage criterion. Uniform \((1-\alpha)\)-coverage means \(\inf_{m\in\mathcal M}\Pr_m\{\tau(m)\in C_n\}\geq1-\alpha\), or an explicitly asymptotic version. The whole parameter space trivially has coverage, so informativeness requires a stated width, risk or power objective. A forecasting improvement is relative to a named benchmark, information set and evaluation population.

\textbf{Equilibrium and optimization.} An equilibrium comprises feasible plans, optimality, beliefs where needed, budget constraints and all market/resource clearing. The primitives or a stated benchmark define each correspondence \(\mathcal E(p;m)\). Nonemptiness is assumed or proved, never inferred just from compact policy sets. A selected-equilibrium target uses a declared selection rule; a robust target quantifies over all permitted equilibria. Use suprema or infima unless attainment is established. An \(\varepsilon\)-optimal policy is feasible and within \(\varepsilon\) of the finite optimum value. Report an empty feasible set separately.

\textbf{Welfare and accounting.} Normative utility, interpersonal normalization, social weights, numeraire, discounting and the use of public receipts are primitives. Behavioral utility need not coincide with the welfare criterion. Household welfare includes ownership income and rents once; adding profits again to compensating variation generally double-counts them. A monetary equivalent is defined by an explicit transfer and price convention, with monotonicity and range conditions. Average effects, mechanism contributions, transfers and welfare losses are different targets. Infinite sums require convergence and dynamic feasible plans require terminal or no-Ponzi conditions.

\textbf{Source versions.} A working-paper date, revision date and journal publication year are distinguished. A DOI for a journal version is not automatically the DOI of an earlier manuscript. Locators refer to the stated version and printed pages where specified. A metadata-only check does not verify a claim about a full-text conclusion. Every entry's source subsection narrows the provenance claim accordingly.
% END ECONOMICS STATEMENT
\end{document}
